\documentclass[../../main.tex]{subfiles}

\begin{document}

\section{逼近与可分性}\label{section:ch07-逼近与可分性}

我们这里详细阐述 Littlewood 第二原理的一般主题, 即用较好的一类函数逼近某类中的函数. 我们考虑关于 $L^{p}(E)$ 范数的逼近. 引入更一般的稠密性的概念是有用的.

\begin{definition}[稠密]\label{definition:ch07-稠密}
令 $X$ 为具有范数 $\|\cdot\|$ 的赋范线性空间. 给定 $X$ 的两个子集 $\mathcal{F}$ 和 $\mathcal{G}$, 满足 $\mathcal{F}\subseteq\mathcal{G}$, 我们说 $\mathcal{F}$ 在 $\mathcal{G}$ 中稠密, 若对 $\mathcal{G}$ 中的每个函数 $g$ 和 $\varepsilon>0$, 存在 $\mathcal{F}$ 中的函数 $f$ 使得 $\|f-g\|<\varepsilon$.
\end{definition}

不难看出, 集合 $\mathcal{F}$ 在 $\mathcal{G}$ 中稠密当且仅当对 $\mathcal{G}$ 中每个 $g$ 存在 $\mathcal{F}$ 中的序列 $\{f_{n}\}$, 使得
\begin{align*}
\text{在 }X\text{ 中 },\quad\lim_{n\to\infty}f_{n}=g.
\end{align*}
此外, 观察到对于 $\mathcal{F}\subseteq\mathcal{G}\subseteq\mathcal{H}\subseteq X$,
\begin{align}
\text{若 }\mathcal{F}\text{ 在 }\mathcal{G}\text{ 中稠密而 }\mathcal{G}\text{ 在 }\mathcal{H}\text{ 中稠密, 则 }\mathcal{F}\text{ 在 }\mathcal{H}\text{ 中稠密}\label{align:ch07-稠密传递}
\end{align}
是有用的.

我们已遇到稠密集: 如同无理数集, 有理数集在 $\mathbf{R}$ 中稠密. 此外, Weierstrass 逼近定理\footnote{关于证明, 见 Patrick Fitzpatrick 的 \emph{Advanced Calculus}[Fit09].}可用当前的赋范线性空间的语言叙述如下: 限制在 $[a,b]$ 上的多项式族在赋予最大值范数的线性空间 $C[a,b]$ 中稠密.

\begin{proposition}\label{proposition:ch07-简单函数稠密}
令 $E$ 为可测集而 $1\leqslant p\leqslant\infty$. 则 $L^{p}(E)$ 中的简单函数构成的子空间在 $L^{p}(E)$ 中稠密.
\end{proposition}

\begin{proof}
令 $g$ 属于 $L^{p}(E)$. 首先考虑 $p=\infty$. 存在 $E$ 的测度为零的子集 $E_{0}$ 使得 $g$ 在 $E\sim E_{0}$ 上有界. 我们从简单逼近引理推出存在 $E\sim E_{0}$ 上的简单函数序列, 它在 $E\sim E_{0}$ 上一致收敛于 $g$, 因此关于 $L^{\infty}(E)$ 范数收敛于 $g$. 于是简单函数在 $L^{\infty}(E)$ 中稠密.

现在假定 $1\leqslant p<\infty$. 函数 $g$ 可测, 因此, 根据简单逼近定理, 存在 $E$ 上的简单函数序列 $\{\varphi_{n}\}$ 使得在 $E$ 上逐点 $\{\varphi_{n}\}\to g$ 且
\begin{align*}
\text{在 }E\text{ 上对所有 }n,\quad|\varphi_{n}|\leqslant|g|.
\end{align*}
从积分的比较判别法得出每个 $\varphi_{n}$ 属于 $L^{p}(E)$. 我们宣称在 $L^{p}(E)$ 中 $\{\varphi_{n}\}\to g$. 事实上, 对所有 $n$, 在 $E$ 上,
\begin{align*}
|\varphi_{n}-g|^{p}\leqslant 2^{p}\{|\varphi_{n}|^{p}+|g|^{p}\}\leqslant 2^{p+1}|g|^{p}.
\end{align*}
由于 $|g|^{p}$ 在 $E$ 上可积, 我们从 Lebesgue 控制收敛定理推出在 $L^{p}(E)$ 中 $\{\varphi_{n}\}\to g$.
\end{proof}

\begin{proposition}\label{proposition:ch07-阶梯函数稠密}
令 $[a,b]$ 为闭有界区间而 $1\leqslant p<\infty$. 则 $[a,b]$ 上的阶梯函数构成的子空间在 $L^{p}[a,b]$ 中稠密.
\end{proposition}

\begin{proof}
前一个命题告诉我们简单函数在 $L^{p}[a,b]$ 中稠密. 因此仅须证明阶梯函数在简单函数中关于 $\|\cdot\|_{p}$ 范数稠密. 每个简单函数是可测集的特征函数的线性组合. 因此, 若每个这样的特征函数可在 $\|\cdot\|_{p}$ 范数下被阶梯函数任意逼近, 由于阶梯函数是线性空间, 任何简单函数也能如此. 令 $g=\chi_{A}$, 其中 $A$ 是 $[a,b]$ 的可测子集且令 $\varepsilon>0$. 我们寻求 $[a,b]$ 上的阶梯函数 $f$, 使得 $\|f-g\|_{p}<\varepsilon$. 根据第 2 章的定理 12, 存在开区间的有限不交族 $\{I_{k}\}_{k=1}^{n}$, 使得若我们定义 $\mathcal{U}=\bigcup_{k=1}^{n}I_{k}$, 则对称差 $A\Delta\mathcal{U}=[A\sim\mathcal{U}]\bigcup[\mathcal{U}\sim A]$ 具有性质:
\begin{align*}
m(A\Delta\mathcal{U})<\varepsilon^{p}.
\end{align*}
由于 $\mathcal{U}$ 是开区间的有限不交族的并, $\chi_{\mathcal{U}}$ 是一个阶梯函数. 此外,
\begin{align}
\|\chi_{A}-\chi_{\mathcal{U}}\|_{p}=[m(A\Delta\mathcal{U})]^{1/p}.\label{align:ch07-特征函数逼近}
\end{align}
因此 $\|\chi_{A}-\chi_{\mathcal{U}}\|_{p}<\varepsilon$, 证明完成.
\end{proof}

\begin{definition}[可分]\label{definition:ch07-可分}
赋范线性空间 $X$ 称为可分的, 若存在一个在 $X$ 中稠密的可数子集.
\end{definition}

实数是可分的, 这是由于有理数集是可数稠密子集. 对闭有界区间 $[a,b]$, 赋予最大值范数的 $C[a,b]$ 是可分的, 这是由于我们从 Weierstrass 逼近定理推出, 具有有理系数的多项式是在 $C[a,b]$ 中稠密的可数集.

\begin{theorem}\label{theorem:ch07-Lp可分}
令 $E$ 为可测集而 $1\leqslant p<\infty$. 则赋范线性空间 $L^{p}(E)$ 是可分的.
\end{theorem}

\begin{proof}
令 $[a,b]$ 为闭有界区间, 而 $\mathcal{S}[a,b]$ 为 $[a,b]$ 上的阶梯函数族. 定义 $\mathcal{S}'[a,b]$ 为 $\mathcal{S}[a,b]$ 的子族. 其中的函数 $\psi$ 在 $[a,b]$ 上取有理值且存在 $[a,b]$ 的分划 $P=\{x_{0},\cdots,x_{n}\}$, 使得对 $1\leqslant k\leqslant n$, $\psi$ 在 $(x_{k-1},x_{k})$ 上是常数, 而对 $1\leqslant k\leqslant n-1$, $x_{k}$ 是有理数. 我们从有理数在实数中的稠密性推出, $\mathcal{S}'[a,b]$ 在 $\mathcal{S}[a,b]$ 中关于 $L^{p}(E)$ 范数稠密. 我们把验证 $\mathcal{S}'[a,b]$ 是可数集作为练习. 以下两个包含关系都关于 $L^{p}[a,b]$ 范数稠密:
\begin{align*}
\mathcal{S}'[a,b]\subseteq\mathcal{S}[a,b]\subseteq L^{p}[a,b].
\end{align*}
因此, 根据\eqref{align:ch07-稠密传递}, $\mathcal{S}'[a,b]$ 在 $L^{p}[a,b]$ 中稠密. 对每个自然数 $n$, 定义 $\mathcal{F}_{n}$ 为 $\mathbf{R}$ 上的函数, 它在 $[-n,n]$ 外消失且限制在 $[-n,n]$ 上属于 $\mathcal{S}'[-n,n]$. 定义 $\mathcal{F}=\bigcup_{n\in N}\mathcal{F}_{n}$. 则 $\mathcal{F}$ 是 $L^{p}(\mathbf{R})$ 中的函数的可数族. 根据单调收敛定理,
\begin{align*}
\text{对所有 }f\in L^{p}(\mathbf{R}),\quad\lim_{n\to\infty}\int_{[-n,n]}|f|^{p}=\int_{\mathbf{R}}|f|^{p}.
\end{align*}
因此, 根据每个 $\mathcal{F}_{n}$ 的选取, $\mathcal{F}$ 是在 $L^{p}(\mathbf{R})$ 中稠密的可数函数族. 最后, 令 $E$ 为一般的可测集. 则 $\mathcal{F}$ 中的函数在 $E$ 上的限制族是 $L^{p}(E)$ 的可数稠密子集, 因此 $L^{p}(E)$ 是可分的.
\end{proof}

如同以下例子所表明的, $L^{\infty}(E)$ 一般是不可分的.

\begin{example}\label{example:ch07-Linfty不可分}
令 $[a,b]$ 为非退化的闭有界区间. 我们宣称赋范线性空间 $L^{\infty}[a,b]$ 是不可分的. 为验证这一断言, 我们用反证法. 假定存在在 $L^{\infty}[a,b]$ 中稠密的可数集 $\{f_{n}\}_{n=1}^{\infty}$. 对每个数 $x\in[a,b]$, 选择自然数 $\eta(x)$ 使得
\begin{align*}
\|\chi_{[a,x]}-f_{\eta(x)}\|_{\infty}<1/2.
\end{align*}
观察到
\begin{align*}
\text{若 }a\leqslant x_{1}<x_{2}\leqslant b,\quad\text{则}\|\chi_{[a,x_{1}]}-\chi_{[a,x_{2}]}\|_{\infty}=1.
\end{align*}
因此 $\eta$ 是 $[a,b]$ 映上自然数集的一一映射. 但自然数集是可数的, 而 $[a,b]$ 是不可数的. 我们从这个矛盾得出, $L^{\infty}[a,b]$ 是不可分的.
\end{example}

对于 $\mathbf{R}$ 的可测子集 $E$, 我们记 $C_{c}(E)$ 为 $E$ 上在一有界集外消失的连续实值函数的线性空间. 在上述定理的证明中, 对于 $1\leqslant p<\infty$, 我们给出了 $L^{p}(\mathbf{R})$ 的稠密子集 $\mathcal{F}$, 它具有性质: 对每个 $f\in\mathcal{F}$, 存在闭有界区间 $[a,b]$, 使得 $f$ 在 $[a,b]$ 上的限制是阶梯函数, 而在 $[a,b]$ 外消失. 不难看出, 每个 $f\in\mathcal{F}$ 是连续分段线性且在有界集外消失的函数序列在 $L^{p}(\mathbf{R})$ 中的极限. 定义 $\mathcal{F}'$ 为 $\mathcal{F}$ 中函数的所有这样的逼近序列的并. 则 $\mathcal{F}'$ 在 $L^{p}(\mathbf{R})$ 中稠密. 此外, 对于可测集 $E$, 属于 $\mathcal{F}'$ 的函数在 $E$ 上的限制所组成的族是由 $E$ 上在有界集外消失的连续函数组成的 $L^{p}(E)$ 的稠密子集. 这证明了以下定理.

\begin{theorem}\label{theorem:ch07-Cc稠密}
令 $E$ 为可测集而 $1\leqslant p<\infty$. 则 $C_{c}(E)$ 在 $L^{p}(E)$ 中稠密.
\end{theorem}

\begin{exercise}\label{exercise:ch07-36}
令 $\mathcal{S}$ 为赋范线性空间 $X$ 的子集. 证明 $\mathcal{S}$ 在 $X$ 中稠密当且仅当每个 $g\in X$ 是 $\mathcal{S}$ 中序列的极限.
\end{exercise}

\begin{exercise}\label{exercise:ch07-37}
验证\eqref{align:ch07-稠密传递}.
\end{exercise}

\begin{exercise}\label{exercise:ch07-38}
证明具有有理系数的多项式族是可数的.
\end{exercise}

\begin{exercise}\label{exercise:ch07-39}
令 $E$ 为可测集, $1\leqslant p<\infty$, $q$ 是 $p$ 的共轭, 而 $\mathcal{S}$ 是 $L^{p}(E)$ 的稠密子集. 证明: 若 $g\in L^{p}(E)$ 且对所有 $f\in\mathcal{S}$, $\int_{E}f\cdot g=0$, 则 $g=0$.
\end{exercise}

\begin{exercise}\label{exercise:ch07-40}
验证定理 11 的证明细节.
\end{exercise}

\begin{exercise}\label{exercise:ch07-41}
令 $E$ 为有限测度的可测集而 $1\leqslant p_{1}<p_{2}<\infty$. 考虑赋予范数 $\|\cdot\|_{p_{1}}$ 的线性空间 $L^{p_{2}}(E)$. 该赋范线性空间是 Banach 空间吗?
\end{exercise}

\begin{exercise}\label{exercise:ch07-42}
展示可测集 $E$ 使得 $L^{\infty}(E)$ 是可分的. 证明: 若 $E$ 包含非退化区间, 则 $L^{\infty}(E)$ 是不可分的.
\end{exercise}

\begin{exercise}\label{exercise:ch07-43}
假定 $X$ 是具有范数 $\|\cdot\|$ 的 Banach 空间. 令 $X_{0}$ 为 $X$ 的稠密子空间. 假设 $X_{0}$ 被赋予继承自 $X$ 的范数也是一个 Banach 空间, 证明 $X=X_{0}$.
\end{exercise}

\begin{exercise}\label{exercise:ch07-44}
对于 $1\leqslant p<\infty$, 证明序列空间 $\ell^{p}$ 是可分的. 证明实数集族是不可数的且得出 $\ell^{\infty}$ 是不可分的.
\end{exercise}

\begin{exercise}\label{exercise:ch07-45}
证明定理 12.
\end{exercise}

\begin{exercise}\label{exercise:ch07-46}
证明对于 $1<p<\infty$ 以及任何两个数 $a$ 和 $b$,
\begin{align*}
\left|\mathrm{sgn}(a)\cdot|a|^{1/p}-\mathrm{sgn}(b)\cdot|b|^{1/p}\right|^{p}\leqslant 2^{p}\cdot|a-b|.
\end{align*}
\end{exercise}

\begin{exercise}\label{exercise:ch07-47}
证明对于 $1<p<\infty$ 以及任何两个数 $a$ 和 $b$,
\begin{align*}
\left|\mathrm{sgn}(a)\cdot|a|^{p}-\mathrm{sgn}(b)\cdot|b|^{p}\right|\leqslant p\cdot|a-b|(|a|+|b|)^{p-1}.
\end{align*}
\end{exercise}

\begin{exercise}\label{exercise:ch07-48}
(Mazur) 令 $E$ 为可测集而 $1<p<\infty$. 对于 $L^{1}(E)$ 中的 $f$, 定义 $E$ 上的函数 $\Phi(f)$ 为
\begin{align*}
\Phi(f)(x)=\mathrm{sgn}(f(x))|f(x)|^{1/p}.
\end{align*}
证明 $\Phi(f)$ 属于 $L^{p}(E)$. 此外, 用习题 46 证明对所有 $L^{1}(E)$ 中的 $f$, $g$,
\begin{align*}
\|\Phi(f)-\Phi(g)\|_{p}^{p}\leqslant 2^{p}\cdot\|f-g\|_{1}.
\end{align*}
由此得出 $\Phi$ 是以下意义的从 $L^{1}(E)$ 到 $L^{p}(E)$ 的连续映射: 若在 $L^{1}(E)$ 中 $\{f_{n}\}\to f$, 则在 $L^{p}(E)$ 中 $\{\Phi(f_{n})\}\to\Phi(f)$. 接着证明 $\Phi$ 是一对一的且它的象是 $L^{p}(E)$. 找到逆映射的公式. 用前一个习题得出逆映射 $\Phi^{-1}$ 是从 $L^{p}(E)$ 到 $L^{1}(E)$ 的连续映射.
\end{exercise}

\begin{exercise}\label{exercise:ch07-49}
用前一个习题证明 $L^{1}(E)$ 的可分性蕴涵 $L^{p}(E)$ 的可分性, 其中 $1<p<\infty$.
\end{exercise}

\begin{exercise}\label{exercise:ch07-50}
对闭有非退化区间 $[a,b]$, 证明不存在从 $L^{1}[a,b]$ 映上 $L^{\infty}[a,b]$ 的连续映射 $\Phi$.
\end{exercise}

\begin{exercise}\label{exercise:ch07-51}
用 Lusin 定理证明定理 12.
\end{exercise}

\end{document}